\section*{3. Picard--Lefschetz 公式}

\subsection*{3.1. 预备结果}

\subsubsection*{3.1.1.}
设 $S$ 为 Hensel 特质, $f:X\to S$ 为有限型平坦态射, 纯相对维数为 $n$. 假设特殊纤维 $X_s$ 在一个有限点集 $\Sigma'$ 之外光滑, 且 $X_s$ 在该点集各点处有普通二次奇点.

设 $\Sigma$ 为那些一般纤维 $X_\eta$ 在 $x$ 的一个邻域内光滑的点 $x\in\Sigma'$ 所成的集合. 最后, 设 $\Lambda$ 为阶与 $S$ 的剩余特征互素的挠环.

\subsubsection*{命题 3.1.2.}
\begin{enumerate}
\item[(i)] \underline{对 $i\ne n$, 层 $R^i\Phi(\Lambda)$ 为零.}

\item[(ii)] \underline{层 $R^n\Phi(\Lambda)$ 在 $\Sigma$ 之外为零. 它在 $\Sigma$ 上的限制为秩一 $\Lambda$-模层.}

\item[(iii)] \underline{对 $\Sigma$ 的每个几何点 $\bar x$, 2.2.5(C) 的配对 $(a,b)$ 使 $\Lambda$-模}
\[
\bigl(R^n\psi(\Lambda)\bigr)_{\bar x}
\quad\text{及}\quad
H^n_{\{\bar x\}}\bigl(X_{\bar s},R\psi_{\bar\eta}(\Lambda)\bigr)
\]
\underline{互为对偶. 若 $n\ne0$, 它们自由且秩为一; 若 $n=0$, 则自由且秩为二.}
\end{enumerate}
\footnote{印本称 $R^n\Phi(\Lambda)$ 在 $\Sigma$ 之外“独有”(\emph{seul}); 证明及其支集在 $\Sigma$ 上的断言要求应为“零”(\emph{nul}). 在 (iii) 中印本写为 $H^n_{\{\bar x\}}(R^n\psi_{\bar\eta}(\Lambda))$; 所引用的 2.2.5(C) 中配对的第二项为 $H^n_{\{\bar x\}}(X_{\bar s},R\psi_{\bar\eta}(\Lambda))$.}

3.1.2 的各断言在平展拓扑下是局部的. 因此只需处理以下情形: $S$ 严格 Hensel, $X$ 为相对仿射空间 $\mathbb A^{n+1}_S$ 中由满足 2.2.1(a) 和 (b) 的二次方程定义的超曲面(应用 1.3.2).

态射 $f$ 在 $X_s-\Sigma'$ 的每点处光滑, 故在 $\Sigma'$ 之外有 $R\Phi=0$. 由 2.2.4, $R\Phi$ 甚至在 $\Sigma$ 之外也为零. 现在可假设 2.2.5(*); 此时断言 (ii) 和 (iii) 由 2.2.5(A)、(B) 和 (C) 得到.